Con los siguientes cambios de variable

\begin{equation*}
    e=\begin{bmatrix}
        \bar{e} \\ e_n
    \end{bmatrix} \quad , \quad \bar{e}=\begin{bmatrix}
        e_1 \\ e_2 \\ ... \\ e_{n-1}
    \end{bmatrix}
\end{equation*}

\begin{equation*}
    F=\begin{bmatrix}
        \bar{f} \\ f_n
    \end{bmatrix}  \quad , \quad \bar{f}=\begin{bmatrix}
        f_1 \\ f_2 \\ ... \\ f_{n-1}
    \end{bmatrix}
\end{equation*}

\begin{equation*}
    m=\begin{bmatrix}
        \bar{m} \\ 1
    \end{bmatrix}  \quad , \quad \bar{m}=\begin{bmatrix}
        m_1 \\ m_2 \\ ... \\ m_{n-1}
    \end{bmatrix}
\end{equation*}

Obtenemos

\begin{equation*}
    \begin{aligned}
        \sigma=e_n+\bar{m}^T\bar{e}=0\\
        e_n=-\bar{m}^T\bar{e}
    \end{aligned}
\end{equation*}

derivando el error

\begin{equation*}
    \dot{\bar{e}}=A_m\bar{e}-\bar{f}
\end{equation*}

con

\begin{equation*}
    A_m=\bar{A}-\bar{B}\bar{m}^T\begin{bmatrix}
        0 & 1 & 0 & \dots & 0 \\
        0 & 0 & 1 & \dots & 0 \\
        \dots & \dots & \dots & \dots &\dots \\
        0 & 0 & 0 & \dots & 1 \\
        -m_1 & -m_2 & m_3 & \dots & -m_{n-1}
    \end{bmatrix}
\end{equation*}

Se elige la función de Lyapunov

\begin{equation*}
    \bar{V}=\bar{e}^T P \bar{e} \quad , \quad P=P^T>0
\end{equation*}

derivando

\begin{equation*}
    \begin{aligned}
        \dot{\bar{V}}&=\bar{e}^TP\dot{\bar{e}}+\dot{\bar{e}}^TPe\\
         &=e^TP[A_m\bar{e}-\bar{F}]+[A_m\bar{e}-\bar{F}]^TP\bar{e}\\
         &=e^T[PA_m+A_m^TP]\bar{e}-2\bar{e}^TP\bar{F}
    \end{aligned}
\end{equation*}

con la desigualdad de Young

\begin{equation*}
    -2\bar{e}^TP\bar{F} \leq \frac{1}{\gamma}\bar{e}^TPP\bar{e}+\gamma||\bar{F}||^2
\end{equation*}

para mantener la condicion de la elipsoide tenemos que $\dot{V}\leq -\alpha V+\beta$

\begin{equation*}
    \dot{\bar{V}}\leq \bar{e}^T[PA_m+A_m^TP-P(-\gamma)^{-1}P]\bar{e}+\gamma||F||^2
\end{equation*}

esta condicion se traduce a

\begin{equation*}
    \bar{e}^T[PA_m+A_m^TP-P(-\gamma)^{-1}P]\bar{e} \leq - \alpha\bar{e}^-{-1}P\bar{e}
\end{equation*}

\begin{equation*}
    PA_m+A_m^TP-P(-\gamma)^-1P+\alpha P \leq 0
\end{equation*}

Proponemos el cambio de variable

\begin{equation*}
    \begin{aligned}
        QP^{-1} >0 \\
        H = A_m Q \\
        G = \bar{m}^T Q
    \end{aligned}
\end{equation*}

tenemos que

\begin{equation*}
    \begin{aligned}
        A_m= HQ^{-1} \\
        Am^T=(Q^{-1})^TH^T\\
        P(-\gamma I)^{-1}P = (Q^{-1})^T(-\gamma I^{-1})Q^{-1}
    \end{aligned}
\end{equation*}

Así mismo

\begin{equation*}
    \begin{aligned}
        \alpha P &= \alpha Q^{-1}I \\
            &= \alpha Q^{-1}QQ^{-1}\\
            &= \alpha (Q^{-1})^TQQ^{-1}
    \end{aligned}
\end{equation*}

con el cambio de variable

\begin{equation*}
    \begin{aligned}
        (Q^-1)^THQ^{-1}+(Q^{-1})^TH^TQ^{-1} - (Q^{-1})^T(-\gamma I)^{-1}Q^{-1}+\alpha Q^{-1} \leq 0 \\
        (Q^{-1})^THQ^{-1}+(Q^{-1})^TH^TQ{-1}-(Q^{-1})^TH(-\gamma I)^{-1}+\alpha (Q^{-1})^TQQ^{-1} \leq 0 \\
        (Q^{-1})^T [H+H^T + \alpha Q - (- \gamma I )^{-1}]Q^{-1} \leq 0 \\
        (Q^{-1})^T [H + H^T + \gamma Q - I(- \gamma I)^{-1} I] Q^{-1} \leq 0
    \end{aligned}
\end{equation*}

con el complemento de schur

\begin{equation*}
    \begin{bmatrix}
        H+H^T+\alpha Q & I \\ I & -\gamma I
    \end{bmatrix} < 0
\end{equation*}

\begin{equation*}
    Q>0
\end{equation*}




